

\begin{unnumbered}

\subsection{1. НАЙМЕНУВАННЯ ТА ОБЛАСТЬ ЗАСТОСУВАННЯ}

У ході HJ дипломного проєкту було проведено аналіз проблематики мережевої взаємодії компонентів інформаційної системи вищого навчальногго закладу.

Було спроектовано і розроблено протокол виклику віддалених процедур та автоматичні тести, що перевіряють його відповідність вимогам. Розроблювана система додає ряд переваг, які не було реалізовано в проаналізованих аналогах, а саме:

Можливість динамічної інтроспекції віддалених інтерфейсів;

Прозорий для користувача спосіб виклику віддалених функцій, що не відрізняється від виклику локальних асинхронних функцій. 

Було проведено аналіз якості програмного забезпечення. Було протестовано розроблений застосунок. Тестування продемонструвало, що протокол реалізовано коректно, розроблене програмне забезпечення відповідає усім функціональним вимогам та готове до розгортання в робочому середовищі

\subsection{2. ПІДСТАВИ ДЛЯ РОЗРОБКИ}

Можливість динамічної інтроспекції віддалених інтерфейсів;

Прозорий для користувача спосіб виклику віддалених функцій, що не відрізняється від виклику локальних асинхронних функцій. 

Було проведено аналіз якості програмного забезпечення. Було протестовано розроблений застосунок. Тестування продемонструвало, що протокол реалізовано коректно, розроблене програмне забезпечення відповідає усім функціональним вимогам та готове до розгортання в робочому середовищі

%\section{2. ПІДСТАВИ ДЛЯ РОЗРОБКИ}
%
%Було спроектовано і розроблено протокол виклику віддалених процедур та автоматичні тести, що перевіряють його відповідність вимогам. Розроблювана система додає ряд переваг, які не було реалізовано в проаналізованих аналогах, а саме:
%
%Можливість динамічної інтроспекції віддалених інтерфейсів;
%
%Прозорий для користувача спосіб виклику віддалених функцій, що не відрізняється від виклику локальних асинхронних функцій. 
%
%Було проведено аналіз якості програмного забезпечення. Було протестовано розроблений застосунок. Тестування продемонструвало, що протокол реалізовано коректно, розроблене програмне забезпечення відповідає усім функціональним вимогам та готове до розгортання в робочому середовищі

\subsection{3. МЕТА ТА ПРИЗНАЧЕННЯ РОЗРОБКИ}

Було спроектовано і розроблено протокол виклику віддалених процедур та автоматичні тести, що перевіряють його відповідність вимогам. Розроблювана система додає ряд переваг, які не було реалізовано в проаналізованих аналогах, а саме:

Можливість динамічної інтроспекції віддалених інтерфейсів;

Прозорий для користувача спосіб виклику віддалених функцій, що не відрізняється від виклику локальних асинхронних функцій. 

Було проведено аналіз якості програмного забезпечення. Було протестовано розроблений застосунок. Тестування продемонструвало, що протокол реалізовано коректно, розроблене програмне забезпечення відповідає усім функціональним вимогам та готове до розгортання в робочому середовищі


\subsection{4. ДЖЕРЕЛА РОЗРОБКИ}

\subsection{5. ТЕХНІЧНІ ВИМОГИ}
\subsubsection{5.1. Вимоги до розробленого продукту}
\subsubsection{5.2. Вимоги до розробленого продукту}
\subsubsection{5.3. Вимоги до програмного забезпечення}
\subsubsection{5.4. Вимоги до апаратної частини}

\subsection{6. ЕТАПИ РОЗРОБКИ}

\end{unnumbered}
